% =====================================================================
%  Chapter 19 — The Calderón–Zygmund logarithm
% =====================================================================
\chapter{\texorpdfstring{The Calder\'on--Zygmund Logarithm}{The Calderon-Zygmund Logarithm}}
\label{ch:cz-log}

\begin{record}
\textbf{The single symbol this book is about.}
\[
  \Lfac \;=\; \log\!\left(e + \frac{\norm{u}_{H^3}}{M}\right),
  \qquad M=\norm{\omega}_{L^\infty}.
\]
It is defined at line~5712 of the proof document, inside
\texttt{lem:local-enstrophy-typeI} (line~5701), and it arrives with the
inequality
\begin{equation}\label{eq:cz-endpoint}
  \norm{\nabla u}_{L^\infty} \;\lesssim\; M\,\Lfac ,
\end{equation}
which is the \(L^\infty\) endpoint of the Biot--Savart operator. Every
enstrophy budget the programme writes from \Sn{33} onward carries it. Three
separate closure attempts fail by exactly one factor of it. It is the audit's
eighth defect, and it is Wall~1.

\smallskip
This chapter is its biography: where it comes from, why it is not removable by
care in the downstream algebra, what kind of object it is in the language of
the renormalisation group, and what \Sn{48} did and did not do to it.
\end{record}

\section{Where it comes from}
\label{sec:cz-origin}

The programme's estimates all have the same shape. A budget is written for a
quantity built from the vorticity --- local enstrophy, angular energy, a
coupling density. The budget's dangerous term is vortex stretching, which
involves \(\nabla u\). And \(\nabla u\) is not available: only \(\omega\) is,
and \(\nabla u\) must be recovered from it by Biot--Savart, which is a
Calder\'on--Zygmund operator.

Calder\'on--Zygmund operators are bounded on \(L^p\) for \(1<p<\infty\). They
are \emph{not} bounded on \(L^\infty\); the correct endpoint is
\(L^\infty\to\mathrm{BMO}\). Recovering a genuine sup-norm therefore costs
something, and what it costs is a logarithm.

\subsection{The mechanism, in one paragraph}

The cleanest derivation in the document is not in \S33 where the factor first
appears, but in \Sn{48}'s \texttt{lem:s48-biotsavart} (line~10047), which
proves a \emph{different} estimate and notices the logarithm's origin on the
way past. Decompose \(v\) into Littlewood--Paley blocks. For the velocity
itself, divergence-freeness gives
\(\Delta_j v = 2^{-j}\mathcal{K}_j\widetilde\Delta_j\omega\), so
\(\norm{\Delta_j v}_{L^\infty}\le C2^{-j}\norm{\omega}_{L^\infty}\) --- and
the geometric factor \(2^{-j}\) makes the sum over blocks converge. For the
\emph{gradient}, the same decomposition gives
\[
  \norm{\Delta_j\nabla v}_{L^\infty} \;\le\; C\norm{\omega}_{L^\infty}
  \qquad\text{with \emph{no} decaying factor,}
\]
and summing over the \(O(\Lfac)\) active blocks produces exactly
\eqref{eq:cz-endpoint}.

\begin{obsv}[The logarithm counts scales]
\label{obs:log-counts-scales}
That is the whole of it, and it is worth saying in words. Each dyadic scale
contributes an equal amount to \(\norm{\nabla u}_{L^\infty}\). Nothing
suppresses the small scales and nothing suppresses the large ones. So the
total is the per-scale contribution times \emph{the number of active scales},
and the number of active scales between the energy scale and the vorticity
scale is \(\Lfac\).

\smallskip
This is why no rearrangement of the downstream algebra removes it: the factor
is not an inefficiency in an estimate, it is a count of how many scales are
simultaneously in play. To remove it one has to change what is being summed,
not how carefully.
\end{obsv}

\section{The genealogy}
\label{sec:cz-genealogy}

The factor was first written down as a criticism, not as a symbol.
Chapter~\ref{ch:inheritance} tells that story in full and
Table~\ref{tab:cz-life} there traces it session by session; this section gives
the mathematical spine.

\begin{figure}[ht]
\centering
\begin{tikzpicture}[
  >={Latex[length=2mm]},
  ev/.style={draw, rounded corners=2pt, font=\scriptsize, align=left,
             inner sep=4pt, text width=52mm},
  bad/.style={ev, draw=impossiblered, thick},
  neu/.style={ev, draw=rulegrey},
  gd/.style={ev, draw=provedgreen, very thick},
  ar/.style={->, thick, rulegrey}
]
\node[bad] (d8) at (0,0)
  {\Sn{31} \textbf{Defect D8}, \texttt{sec:audit-s31}\\
   ``\(\norm{\nabla u}_{L^\infty}\le CM\) \dots\ the correct bound carries a
   logarithmic factor''};
\node[neu, below=5mm of d8] (marg)
  {\Sn{32} \texttt{rem:marginality} (5443)\\
   the logarithm is \emph{mandatory}; the direction equation is exactly
   critical at \(\rstar\) up to it};
\node[neu, below=5mm of marg] (born)
  {\Sn{33} \texttt{lem:local-enstrophy-typeI} (5701)\\
   \(\Lfac\) defined at 5712; enters the machinery};
\node[bad, below=5mm of born] (crude)
  {\Sn{34} \texttt{prop:crude-K} (5903)\\
   \emph{first rediscovery}: misses closing the loop by exactly one logarithm};
\node[bad, below=5mm of crude] (ann)
  {\Sn{41} \texttt{rem:s41-shortfall} (8428)\\
   \emph{second}: absorption needs \(C\Lfac\le\delta\nu\); fails as
   \(\Lfac\to\infty\)};
\node[bad, below=5mm of ann] (band)
  {\Sn{44} \texttt{prop:s44-linfty-route} (9242)\\
   \emph{third}: closure needs \(c_K\le C_1/\Lfac\) --- \emph{decay}, not
   boundedness};
\node[neu, below=5mm of band] (meas)
  {\Sn{45} \texttt{scale\_diagnostics.py}\\
   slope of \(\log c_K\) on \(\log\Lfac\): \(-0.196\), \(+0.088\);
   required \(-1\). Flat};
\node[neu, below=5mm of meas] (w1)
  {\Sn{47} \textbf{Wall 1} --- \OPEN{}, structural.\\
   One obstacle seen three times, not three obstacles};

\foreach \a/\b in {d8/marg, marg/born, born/crude, crude/ann, ann/band,
                   band/meas, meas/w1}{\draw[ar] (\a) -- (\b);}

\node[gd, right=9mm of crude, text width=48mm] (s48)
  {\Sn{48} \texttt{thm:s48-R}\\
   On the \emph{profile} branch the logarithm is never \emph{incurred}: the
   route does not pass through the CZ operator. Wall~2 \WITHDRAWN{}};
\node[bad, below=6mm of s48, text width=48mm] (s48b)
  {\dots\ but \texttt{obs:s48-wall1} (10400): \S33--\S37 need \(\nabla u\)
   against \(M(t)\), not \((T-t)^{-1}\). \textbf{Wall 1 is untouched}};
\draw[->, thick, provedgreen] (band.east) to[out=20,in=-90] (s48.south);
\draw[->, thick, impossiblered, dashed] (s48b.west) to[out=180,in=0] (w1.east);
\end{tikzpicture}
\caption[The life of the logarithm]{From a clause in a defect list to the
programme's central structural barrier, and the one branch on which it turned
out never to have been incurred. The dashed red arrow is the point most easily
lost: \Sn{48}'s result does not reach Wall~1.}
\label{fig:cz-genealogy}
\end{figure}

\subsection{Why the three appearances are one}

The three closure attempts differ completely in their subject matter. The
first is about absorbing a coupling term in a Caccioppoli inequality for the
angular energy. The second is about a residue of level-set truncation on a
band of vorticity magnitudes. The third is about whether a correlation
constant on the transition band is small.

They fail identically. In each case, the quantity that the second rung is able
to absorb is \(\delta\nu(\rstar)^3\norm{w}_{L^\infty}\), and the quantity the
budget delivers is \(C\Lfac\norm{w}_{L^\infty}(\rstar)^3\). Absorption
requires \(C\Lfac\le\delta\nu\), which fails for fixed \(\nu\) as
\(\Lfac\to\infty\). \texttt{prop:s41-crude} ``reproduces, term for term, the
bound of \texttt{prop:crude-K}''.

That identity is the substance of Wall~1's classification. The wall does not
say the logarithm is unbeatable; it says that beating it in one of the three
places means beating it in all three, and that treating them as separate
problems is an error.

\subsection{Decay, not boundedness}

The sharpest form of the obstruction is \Sn{44}'s, and it is worth isolating
because it is what makes the numerical programme of
Chapter~\ref{ch:numerical-lessons} futile in principle rather than merely in
practice.

\texttt{prop:s44-linfty-route} (line~9242) closes the band route if and only
if \(c_K^{\mathrm{band}}\le C_1/\Lfac\). Not \(O(1)\): \emph{decaying}. The
reason is visible in the proof's arithmetic. The region measure cancels
exactly; the enstrophy budget contributes its \(\Lfac M^{1/2}\); and the
hypothesised \(C_1/\Lfac\) is what cancels the \(\Lfac\) the budget carries,
landing the term in the class the second rung absorbs. The hypothesis is being
asked to do exactly one job, and that job is to eat the logarithm. A bounded
correlation constant does not eat it. It leaves it exactly where it was.

\section{A marginal logarithm}
\label{sec:marginality-rg}

\begin{record}
\textbf{Status of this section.} What follows is a way of thinking, recorded
in the \Sn{46} session log as ``the right mental model'' and reproduced here
because it makes the shape of the difficulty legible. \emph{It is not
mathematics in this programme}: no statement in the proof document depends on
it, and nothing in this book is proved with it. A reader who dislikes physical
analogies can skip to \S\ref{sec:s48-R} without loss.
\end{record}

\subsection{The dimensional accident}

\texttt{rem:marginality} (line~5443) records the calculation. On the
self-similar cylinder \(Q(\rstar)=B_{\rstar}\times[t_0-(\rstar)^2,t_0]\), with
\(\rstar=M^{-1/2}\), the stretching term in the direction equation acts as a
potential of size \(\norm{\nabla u}_{L^\infty}\), and its Gr\"onwall exponent
over the cylinder's time span is
\[
  \norm{\nabla u}_{L^\infty}\cdot(\rstar)^2
  \;\lesssim\; M\cdot M^{-1}\cdot\Lfac \;=\; \Lfac .
\]
The powers of \(M\) cancel \emph{exactly}. What is left is a logarithm, and a
logarithm alone.

That exact cancellation is what ``marginal'' means. A quantity that scaled as
a positive power of \(M\) would be irrelevant --- it would die as the scale
shrinks, and the estimate would close trivially. A quantity that scaled as a
negative power would be relevant, would blow up, and the route would be dead.
The programme's quantity does neither. It sits exactly at the boundary, and
whether it is fatal is decided by what a dimensionless, slowly varying factor
does.

\subsection{The analogy, and its limits}

In the language of the renormalisation group, a coupling whose classical
scaling dimension vanishes is \emph{marginal}, and the signature of
marginality is precisely a logarithm: the coupling does not scale, it
\emph{runs}, and it runs logarithmically in the scale. Three-dimensional
Navier--Stokes at the self-similar scale is a marginal-dimension problem in
exactly this sense --- structurally the same position as a four-dimensional
gauge theory, where the coupling is classically dimensionless and everything
depends on the sign of the logarithmic running.

Marginal theories go two ways.

\begin{itemize}[leftmargin=1.6em]
\item \emph{Marginally irrelevant.} The logarithms resum with a favourable
  sign; the effective coupling weakens as the scale shrinks; the theory is
  controlled in the ultraviolet. This is asymptotic freedom.
\item \emph{Marginally relevant.} The logarithms resum with the other sign;
  the effective coupling strengthens without bound at short distance. This is
  the Landau-pole scenario.
\end{itemize}

\begin{obsv}[What \(c_K\le C/\Lfac\) is, in this reading]
\label{obs:ck-is-resummation}
The programme's remaining closure hypothesis says that a correlation constant
\emph{decays like one over the logarithm}. Read through the analogy above,
that is precisely the statement that the marginal coupling is marginally
\emph{irrelevant}: that the accumulation of equal contributions from each of
the \(\Lfac\) active scales is offset by a decorrelation between those scales
that grows at exactly the same rate.

\smallskip
So the question ``does the marginal coupling resum harmlessly or pile up?''
and the question ``is \(c_K\le C/\Lfac\)?'' are the same question, asked in
two vocabularies. The programme has posed it precisely, and has not answered
it. \Sn{45} measured the closest accessible proxy and found the constant
\emph{flat} in \(\Lfac\) over the range a regular solver can reach --- which
Chapter~\ref{ch:numerical-lessons} explains is exactly what both answers
predict on such a flow, and therefore no evidence at all.
\end{obsv}

The analogy's limits should be stated with the analogy. There is no
Navier--Stokes renormalisation group in the sense there is one for a
quantum field theory: no exact scale transformation with a well-defined
generator, no beta function, and no theorem that the logarithms exponentiate.
What survives the analogy is a way of asking the question and a warning about
what kind of answer would be needed --- an exact cancellation at leading
order, not a better constant.

\section{\texorpdfstring{\Sn{48}: the branch where it is never incurred}{S48: the branch where it is never incurred}}
\label{sec:s48-R}

Chapter~\ref{ch:profile-rigidity}'s \Sn{39} reframe carried, in
\texttt{prop:transfer-audit}(d), the observation that on a Type-I profile
``the logarithmic factor \(\Lfac\) may be replaced by an absolute constant''.
That was labelled \PROVEDMOD{(R)}, and \textbf{(R)} was itself an assumed
log-free gradient bound --- which is why \Sn{47} recorded it as Wall~2: the
barrier relocated, not removed.

\Sn{48} settled \textbf{(R)}.

\subsection{The theorem}

\begin{statement}[\texttt{thm:s48-R}, \S\,\texttt{sec:hypR-s48}]
\label{st:s48R}
\PROVED{}. Assume the standing Type-I bounds of line~5565 --- \emph{both}
\(M(t)\le C_I/(T-t)\) and \(\sup_{\T^3}\abs{u(\cdot,t)}\le C_I(T-t)^{-1/2}\)
--- and let \(u_\lambda\) be the rescaled family. Then for every
\(k,l\ge0\) there is \(C_{k,l}(c_I)\) with
\[
  \norm{\nabla^k\partial_s^{\,l}u_\lambda(\cdot,s)}_{L^\infty(\Lambda_\lambda)}
  \;\le\; C_{k,l}(c_I)\,\abs{s}^{-\frac{k+1}{2}-l},
\]
for every \(\lambda\in(0,1]\) and every admissible \(s<0\). Taking \(k=1\),
\(l=0\) is exactly \textbf{(R)} --- \emph{without} the factor \((1+R)\) and
\emph{without} any dependence on \(S_2\).
\end{statement}

\subsection{Why the named trap never bit}

The hypothesis had been carried for nine sessions, and the reason nobody
discharged it was a specific and correct worry:
\(\varepsilon\)-regularity theorems require \emph{smallness} of a scaled local
energy, while Type-I supplies only \emph{boundedness}. That trap is real. This
route never touches it.

\begin{record}
\textbf{\texttt{rem:s48-subcritical}} (line~9822). The bound actually
available on the rescaled family is an \(L^\infty_tL^\infty_x\) \emph{velocity}
bound, so the exponent pair is \(s'=s=\infty\) and
\[
  \frac{2}{s'}+\frac{3}{s} \;=\; 0 \;<\; 1 .
\]
This is the \emph{strict-inequality} range of Serrin's interior regularity
criterion --- ``the range Serrin himself proved in 1962, not the borderline
equality case that needed Struwe and Takahashi''. Boundedness suffices; no
smallness is required, and none is available. ``This single observation is
what decides the present question, and it is why the \(\varepsilon\)-regularity
literature is the wrong shelf: those theorems live at the critical scale, where
smallness is the whole content.''
\end{record}

The quantitative, scale-invariant form comes from the \(L^\infty\)-based local
theory of \cite[\S4]{KNSS2009} and \cite[Thm.~2.4]{SereginSverak2009}, neither
of which uses a local energy, a pressure norm or a smallness hypothesis. The
mechanism is a restart argument on the KNSS estimate, transferred to
\(\Lambda_\lambda=\lambda^{-1}\T^3\) by a periodised Oseen kernel whose only
estimate, \(\int\abs{K^{(L)}}\le Ct^{-1/2}\), is uniform in the period;
parasitic solutions are excluded outright because a harmonic function on a
closed manifold is constant.

\subsection{The correction to the document's own reasoning}

\begin{record}
\textbf{\texttt{rem:s48-mechanism}} (line~10315). The proof of
\texttt{prop:transfer-audit}(d) had explained the disappearance of \(\Lfac\)
by saying that ``\(U\) has, by construction, no structure below the scale
\(\abs{s}^{1/2}\) that is not already visible at scale \(\abs{s}^{1/2}\)''.
That is a heuristic, it is not proved anywhere in the document, and \emph{it
is not the reason}. The reason is the strict subcriticality above:
\(\nabla U\) is obtained by \emph{parabolic smoothing} from \(U\), not by a
Calder\'on--Zygmund bound on \(\omega_U\), and parabolic smoothing has no
endpoint failure.

\smallskip
The correct one-line statement: \emph{the logarithm is not removed on the
profile; it is never incurred, because the profile route does not pass through
the Calder\'on--Zygmund operator at all.}
\end{record}

The correction strengthens the result and is a good example of the source
hierarchy doing its job: a conclusion that was right, supported by a reason
that was wrong, found by checking the reason rather than the conclusion.

\subsection{Consequences}

\texttt{prop:transfer-audit}(d), first clause: \PROVEDMOD{(R)} becomes
\PROVED{}. \textbf{Wall~2} moves from \OPEN{} to \WITHDRAWN{}. Every row of
the \Sn{39} ledger reading ``mod (R)'' loses the (R)
(Table~\ref{tab:s39-ledger}).

\section{The three things that did not happen}
\label{sec:s48-limits}

This is the part of the \Sn{48} result that a summary would drop, and it is
the part that decides what the programme should do next.

\subsection{Part (e) is proved modulo saturation}

\texttt{prop:transfer-audit}(d) states \emph{two} things. The first,
\(\norm{\nabla U}_{L^\infty}\le C_1(c_I)\abs{s}^{-1}\), is now proved. The
second, \(\norm{\nabla U(\cdot,s)}_{L^\infty}\le C(c_I)M_U(s)\), carries a
\emph{pre-existing} proviso: ``whenever the Type-I upper bound is saturated''.
Part~(e) uses the second form. The two differ by
\(\left(\abs{s}M_U(s)\right)^{-1}\), which is \(O(1)\) only on the saturated
set. So (e) is \PROVEDMOD{saturation}, not \PROVED{}. Whether
\(M_U(s)\gtrsim\abs{s}^{-1}\) holds for all \(s<0\) or only on a set is not
settled by \textbf{(R)} and is not settled anywhere.

\subsection{Wall 1 is untouched}

\begin{record}
\textbf{\texttt{obs:s48-wall1}} (line~10400), recorded in the document and
deliberately used nowhere. The same estimate applies verbatim to \(u\) itself
on \(\T^3\), giving, under the standing hypotheses,
\[
  \norm{\nabla u(\cdot,t)}_{L^\infty(\T^3)} \;\le\; C(c_I)\,(T-t)^{-1},
  \qquad\text{logarithm-free.}
\]
``This does \emph{not} remove Wall~1. Every estimate of \S33--\S37 needs
\(\nabla u\) measured against \(M(t)\), not against \((T-t)^{-1}\), and the two
agree only where the Type-I upper bound is saturated --- the same proviso as
residue~(1) above, one level down.''
\end{record}

\begin{obsv}[Saturation, not the logarithm, may be what both branches need]
\label{obs:saturation-is-the-target}
The two obstructions in the previous two subsections are the same statement at
two levels: on the profile, is \(M_U(s)\gtrsim\abs{s}^{-1}\)? On the torus, is
\(M(t)\gtrsim(T-t)^{-1}\)? Both ask whether the Type-I \emph{upper} bound is
matched by a lower bound of the same order --- whether the singularity is
saturating its own rate.

\smallskip
That question is cheaper and better defined than either wall. It is a
statement about \emph{one scalar function of one variable}. It did not exist
as a named question before \Sn{48}. It is the sharpest live target the
programme has, and Chapter~\ref{ch:where-it-stands} says so.
\end{obsv}

\subsection{The ceiling is unchanged}

\texttt{rem:s39-distance} stands verbatim, including its Type-II reservation.
Wall~3 is unaffected. \texttt{rem:s48-scope} says so directly: the section
``does not close either horn of the pincer, does not bear on
\(\sstar\)-decay, does not touch Type-II, and does not bring the programme
nearer to the Clay problem than \texttt{rem:s39-distance} already states''.

\section{A route permanently closed}
\label{sec:bridge-fails}

One further \Sn{48} result belongs here, because it is about the logarithm and
because it forecloses an attempt a future session would otherwise make.

The programme's Type-I hypothesis is often \emph{stated} on the vorticity, and
it is natural to ask whether the velocity bound at line~5565 could be derived
rather than assumed. It cannot.

\begin{statement}[\texttt{prop:s48-bridge-fails}, line~10086]
\label{st:bridge-fails}
\IMPOSSIBLE{}. If
\(\norm{v}_{L^\infty}\le C\norm{\operatorname{curl}v}_{L^\infty}^{\theta}\norm{v}_{L^2}^{1-\theta}\)
holds for all smooth compactly supported divergence-free \(v\) on \(\R^3\),
then \(\theta=3/5\) --- forced by homogeneity under
\(v_\mu(x)=\mu^a v(\mu x)\), and attained. Consequently, from the vorticity
Type-I bound and the energy bound alone, the best obtainable velocity bound is
\(\norm{u}_{L^\infty}\le CC_I^{3/5}E_0^{1/5}(T-t)^{-3/5}\); and
\(3/5>1/2\), so the Type-I velocity rate is \emph{not} implied. The rescaled
family then satisfies only \(\abs{u_\lambda}\le C\lambda^{-1/5}\abs{s}^{-3/5}\),
which diverges as \(\lambda\downarrow0\): no uniform bound, and no limit
object.
\end{statement}

The irony is worth recording. The velocity bridge \emph{is} logarithm-free ---
\texttt{lem:s48-biotsavart} proves it, and the reason is exactly the
convergent \(2^{-j}\) sum of \S\ref{sec:cz-origin} --- and it is still
useless, because the exponent it forces is the wrong one. The near-field
heuristic (near field \(\lesssim\rstar M\) with \(\rstar=M^{-1/2}\), giving
\(M^{1/2}=(T-t)^{-1/2}\)) is correct as far as it goes; it is the \emph{far}
field that fails, contributing \((T-t)^{-3/4}E_0^{1/2}\), worse than the near
field, and pushing the optimal balance to \(3/5\).

So the velocity Type-I bound is an independent hypothesis and must stay one.

\section{Two miscitations, found on the way}
\label{sec:cz-miscitations}

\texttt{rem:s48-miscitations} (line~10011) records two citation defects in the
document's own text, found while looking for \textbf{(R)} on the wrong shelf.
They are recorded here because they are the kind of error that a book like
this exists to catch.

\begin{itemize}[leftmargin=1.6em]
\item \cite{Seregin2012} is cited at three places as authority for Type-I
  local regularity technology --- lines 6810, 6848 and 7689.\footnote{The
  document's own remark says the paper is cited ``twice'' and then names three
  line numbers; the session log likewise says twice and names two. Three
  citing sites exist. The discrepancy is immaterial to the finding and is
  recorded because this book's rule is to record them.} The paper is a
  \emph{critical-endpoint} \(L^3\) blow-up result --- its abstract states a
  necessary condition \(\lim_{t\uparrow T}\norm{v(\cdot,t)}_{L^3}=\infty\) ---
  and it contains no Type-I-improved \(\varepsilon\)-regularity theorem. The
  intended reference is presumably \cite{SereginSverak2009} or
  \cite{AlbrittonBarker2019}. \textbf{(N)} is unaffected: it remains assumed.
\item \cite{LadyzhenskayaSeregin1999} is offered at line~6798 as a
  ``sharpening'' of Serrin's criterion by local-pressure theory. It is an
  \(\varepsilon\)-regularity paper for \emph{suitable} weak solutions requiring
  \(p\in L^{3/2}_{\mathrm{loc}}\) and a smallness hypothesis. It is not the
  tool, and no local-pressure theory is needed at all.
\end{itemize}

A third item in the same family: \textbf{(R)}'s own caveat~(ii) --- that the
\(\lambda\)-uniformity must be checked against \(u_\lambda\) having no uniform
global energy bound --- is \emph{void}. The premise is true and the inference
invalid: the \(L^\infty\)-based local theory uses no energy anywhere, global
or local.

\begin{obsv}[The shape of the whole chapter]
\label{obs:cz-shape}
The Calder\'on--Zygmund logarithm was identified as a defect in \Sn{31},
formalised as a symbol in \Sn{33}, met as an apparently new difficulty three
separate times, consolidated as a wall in \Sn{47}, and shown in \Sn{48} never
to have existed on one of the two branches --- for a reason that had nothing
to do with any of the sixteen sessions of work aimed at it, and everything to
do with a criterion Serrin proved in 1962.

\smallskip
On the other branch it is exactly where it was. What changed at \Sn{48} is not
the size of the obstacle but the identity of the question: from ``can the
logarithm be removed?'' to ``is the Type-I upper bound saturated?''. Whether
that is progress will be decided by whether the new question is easier. Nobody
knows yet, and this book does not claim otherwise.
\end{obsv}
